\begin{titlepage}\centering
{\Large\bfseries UNIVERSIDAD NACIONAL DEL ALTIPLANO\par}
\vspace{3mm}
{\normalsize FACULTAD DE INGENIERÍA MECÁNICA ELÉCTRICA,\\ELECTRÓNICA Y SISTEMAS\par}
{\normalsize ESCUELA PROFESIONAL DE INGENIERÍA MECÁNICA ELÉCTRICA\par}
\vspace{7mm}\includegraphics[height=27mm]{figuras/escudo-unap.png}\par
\vspace{7mm}{\large\bfseries CENTRALES HIDROELÉCTRICAS\par}
\textcolor{azul}{\rule{\textwidth}{1pt}}\par\vspace{5mm}
{\LARGE\bfseries\color{azul} PROYECTO ACADÉMICO DE DISEÑO\\DE UNA MINICENTRAL HIDROELÉCTRICA EN LIMBANI\par}
\vspace{5mm}
{\large Replanteamiento del aprovechamiento y suministro aislado\\a Limbani y Huancasayani\par}
\vspace{4mm}
{\normalsize Primera unidad: recurso hídrico, demanda, dimensionamiento,\\regulación, seguridad y evaluación de alternativas\par}
\vspace{8mm}
\begin{tabular}{L{4.1cm}L{9cm}}
\textbf{Estudiante} & Aquiles Taylor Ramos Yapo\\
\textbf{Código de estudiante} & 227368\\
\textbf{Docente} & Mg. Ing. Gregorio Meza Marocho\\
\textbf{Sección} & GRUPO A\\
\textbf{Ubicación} & Limbani, Sandia, Puno, Perú\\
\textbf{Revisión} & R03 - 14 de octubre de 2026\\
\textbf{Elaboración} & Trabajo académico estudiantil\\
\end{tabular}
\vfill\textcolor{azul}{\rule{\textwidth}{.5pt}}\par
{\small PUNO - PERÚ | 2026\par}
\end{titlepage}
